\section[Preposition]{Preposition (Verhältniswort)} \label{sec:deprep}
The versatility \bt{Deutsch} offers makes it impossible to list out all  prepositions, so only the main concepts along with the most frequently used candidates will be touched upon. 

First, we look at \bt{her/hin} correspondence:
\begin{enumerate}
    \item \bt{her}: Movement towards the speaker
    \item \bt{hin}: Movement away from the speaker
\end{enumerate}
These prefixes can be slammed upon prepositions to create new ones with further specification:
\begin{example}
    \begin{tabular}{c c|c c}
        \bt{Heraus} & Out towards me & \bt{Hinaus} & Out away from me (beyond)\\
        \bt{Herein} & In towards me & \bt{Hinein} & In away from me
    \end{tabular}
\end{example}
The word \bt{herein} is consequently used for inviting someone into the room you are currently in. This difference is subtle, and rarely captured in a translator. Since these are also prepositions, they can be coupled with \bt{trennbar} verbs:
\begin{example}
    \bt{herauf/beschwören}: To conjure up (towards me)
\end{example}
For some place prepositions like `behind, under\ldots', there are different forms based on whether they are used as standalones or as part of a noun unit:
\begin{example}
    \begin{tabular}{c c|c c}
        \textit{Noun unit} & \textit{Alone} & \textit{Lone usage} & \textit{Unit usage}\\\hline
        hinter & hinten & Ich stehe hinten & Ich stehe hinter dem Haus\\\hline
        unter & unten & Die Katze ist unten& Die Katze ist unter der Treppe
    \end{tabular}
\end{example}
Standalones are used when a general location  is specified: in other words when case $7$ is used without an associated noun. For example: \bt{ich stehe hinten} just means that I am standing somewhere behind, but a relative placement with respect to other noun landmarks is not specified. In this respect, one may also think of standalone prepositions as indefinite pronoun contractions: \bt{hinter etwas/hinten} 

We now consider \bt{da}-composites, which are a shorthand way to combine a preposition and \bt{das} ($\vm{P}[c,In]$) as a prefix in the form  \bt{da}. To form them, grab \bt{da} and a preposition of choice, and check if they chain a vowel (have consecutive vowels) in between:
\begin{itemize}
    \item \textit{Yes} - add an \bt{r} in between and smash them together
    \item \textit{No} - smash them together
\end{itemize}
\begin{example}
    \begin{tabular}{c c|c c|c c}
        \bt{dahin*} & towards there &\bt{darauf} & on that & \bt{dazu} & to that \\\hline \bt{daher*} & from there &\bt{damit} & with that& \bt{daraus} & from that
    \end{tabular}
\end{example}
Similarly, one may also combine \bt{was} in the form  \bt{wo} as a prefix with prepositions:
\begin{example}
    \begin{tabular}{c c|c c|c c}
        \bt{wohin*}& where to& \bt{worauf} & on what & \bt{wozu} & for what \\\hline \bt{woher*}& where from& \bt{womit} & with what & \bt{woraus}& from what
    \end{tabular}
\end{example}
You might notice that \bt{wozu} doesn't translate as expected. The truth is, although the given meanings are valid for such words, they are often not their only meanings. \bt{Damit} can also mean `so that', and \bt{dazu} can mean `for that purpose'. These double meanings have no pattern and must be learned over the course of time as is. Other combination variants not following the above rule like \bt{somit, infolge, zugunsten} also exist, sometimes with their own double meanings.

The entries with the asterisk (*) represent direct composition, which is a unique feature related to \bt{hin/her}. The \bt{da} or \bt{wo} in these entries are literal, and do not represent a contraction of \bt{das} or \bt{was}. Hence, their meanings deviate from  other examples, despite looking similar.

According to their meanings, different prepositions require  their associated nouns to be in certain cases. Often,  prepositions accept multiple cases, but some accept only one or two, which in rare cases may lead to forced casing contrary to expectation. First, a list of common prepositions:
\begin{setdef}
    \begin{tabular}{c|m{4cm}|c|m{5cm}}
        \textit{Verhältniswort}& \textit{Meaning}& \textit{Cases}& \textit{Sentence}\\\hline
        auf & \bt{on} & $2, 7$ & Ich klettere auf den Berg\\\hline
        von & \bt{from (location); by (passive)} & $5; 3$ &Der Zug fährt von Gleis drei ab\\\hline
        aus & \bt{from (boundary)} & $5$ & Sie kommt aus Deutschland\\\hline
        vor & \bt{before, in front of} & $2, 7$ & Er steht vor mir\\\hline
        in & \bt{in} & $2, 7$ & Wir gehen in die  Höhle\\\hline
        an & \bt{on, at} & $2, 7$ & Am(an + dem) Bahnhof\\\hline
        bei & \bt{at, in case of} &  $7$ & Sie blieb bei seiner Seite\\\hline
        nach & \bt{towards} & $4$ & Wir fahren nach Garmisch\\\hline
        zu & \bt{to, on, at} & $4, 7$ & Wir gehen zu Fuß\\\hline
        mit & \bt{with} & $3, 4$ & Ich komme mit\\\hline
        seit & \bt{since} & $5$ & Seit diesem Zeitalter\ldots\\\hline
    \end{tabular}\end{setdef}\begin{setdef}
    \begin{tabular}{c|m{4cm}|c|m{5.5cm}}
        für* & \bt{for} & $4[2]$ & Für das Reich!\\\hline
        %zugunsten & \bt{in favour of} & $6$ & Deren zugunsten\ldots\\\hline
        um* & \bt{around, at(time)} & $7[2]$& Um den Hals\ldots\\\hline
        mittels & \bt{by means of}& $6$& Mittels des Fahrrads\ldots\\\hline
        bezüglich & \bt{in consideration of (regarding)}& $6$& Bezüglich des Makels\ldots\\\hline
        jenseits & \bt{beyond the boundary of} & $6$ & Jenseits der Stadtgrenze\ldots
    \end{tabular}
\end{setdef}
Anytime the corresponding meaning needs to be described with an `of' at the end, case $6$ must be used. Hence, \bt{mittels} needs case $6$ despite technically being closer to case $3$. The ones marked with an asterisk (*) are exceptions, where the meaning and the [forced case] do not match. 

\bt{Aus} implies an exit from an enclosed boundary, whereas \bt{von} implies an exit from a certain starting point (like a person). \bt{Von} is also used as a case $1$ marker for a passive sentence, and a case $6$ marker for noun units without any additional features. For the second use, the case $6$ noun stays in \bt{Dativ}, since \bt{von} does not support \bt{Genitiv}. \bt{Aus} can also be used when describing the material composition of any substance: `made (out of) glass'

It is possible for \bt{zu} to be used as an emphasizer rather than a preposition: \bt{zu viel Energie} (too much energy). Here, \bt{Energie} is not in case $4$. Such uses are most likely impossible for any  preposition other than \bt{zu}, since I have never come across any other examples.

Finally, it is sometimes possible to place prepositions at the end of a noun unit. This is more common in specific fixed expressions and not generally done:
\begin{example}
    meiner Meinung nach, den Bach herunter\ldots
\end{example}